\documentclass[12pt]{article}
\usepackage[utf8]{inputenc}
\usepackage[T1]{fontenc}
\usepackage{hyperref}
\title{Yaka Language Resources for K-12 Curriculum Development}
\author{Compiled by Claude Sonnet 5.5 (claude-sonnet-5-5) for Andr\'as Kornai}
\date{}
\begin{document}
\maketitle

\section{Introduction}
Yaka (ISO 639-3 yaf; also Iyaka, Iaca, Kiyaka) is a Bantu language of southwestern DR Congo and northern Angola with about 900,000 speakers. Omniglot gives two dialects, Yaka proper and Ngoongo. This report is based on a few web searches and found little. No school materials were found. Searches for ``Kyaka'' also return an unrelated language of Papua New Guinea.

\section{Primary Educational Resources}
No primers, textbooks or school programmes for Yaka were found in the search results.

\section{Linguistic and Reference Materials}
\begin{itemize}
\item Omniglot, Yaka (Yiy\'ak\'a): an overview of the language, its writing and its sound system. \url{https://omniglot.com/writing/yaka.htm}
\item Omniglot chart, Yaka alphabet and tones (PDF). \url{https://www.omniglot.com/charts/yaka.pdf}
\item Glottolog record for Kyota, Kutumisa B. (1999), a morphology and tone study of Yaka (Congo and Angola) listed under the Yaka headings. The result list gave only the record, so the work's content could not be confirmed. \url{https://glosbe.com/} is not cited. Record: \url{https://glottolog.org/resource/reference/id/127261}
\item Webonary Yaka--French dictionary, with alphabet, entries guide and foreword. Caution: it describes Yaka as spoken in the Lekoumou Department, Sibiti District, Republic of Congo. That may be a different variety from the DRC/Angola Yaka, so check before using it. \url{https://www.webonary.work/yaka/overview/alphabet/?lang=en}, \url{https://www.webonary.work/yaka/overview/entries-explained/}, \url{https://www.webonary.work/yaka/overview/foreword/?lang=en}
\end{itemize}

\section{Religious and Community Resources}
None were found in the search results.

\section{Digital and Audio Resources}
No audio, video, corpus or software resources were found in the search results.

\section{Curriculum Development Context}
The searches turned up only descriptive and lexical material, mostly in French or about a possibly different variety. Mother-tongue K-12 materials would need new work, starting with confirming the DRC/Angola orthography and comparing the Webonary lexicon with the local variety. Further searching is needed in French, Portuguese and Kikongo-area sources.
\end{document}